\documentclass[11pt,a4paper]{article}
\usepackage[T1]{fontenc}
\usepackage[utf8]{inputenc}
\usepackage{lmodern,amsmath,amssymb}
\usepackage[margin=25mm]{geometry}
\usepackage[hidelinks]{hyperref}
\hypersetup{pdftitle={Newton's inscribed polygon differs from the parabola by the phase of its segments},pdfauthor={navstokgap research working notes}}
\newcommand{\tightlist}{\setlength{\itemsep}{0pt}\setlength{\parskip}{0pt}}
\setlength{\emergencystretch}{3em}
\setcounter{secnumdepth}{0}
\begin{document}
\hypertarget{newtons-inscribed-polygon-differs-from-the-parabola-by-the-phase-of-its-segments}{%
\section{Newton's inscribed polygon differs from the parabola by the
phase of its
segments}\label{newtons-inscribed-polygon-differs-from-the-parabola-by-the-phase-of-its-segments}}

For a body of mass \(m\) under a constant transverse force \(F\) on
\([0,\tau]\), take any partition into steps \(\tau_1,\dots,\tau_N\) and
Newton's polygon of Proposition I inscribed in the Galileo parabola:
uniform motion along each chord, impulses at the vertices. Classically
the polygon and the curve are in the same state at every vertex. Quantum
mechanically their evolutions differ by a pure phase, the same for every
state of the body:

\[\boxed{W_{\rm poly}=e^{i\theta_N}\,W_F,\qquad
\theta_N=\frac{F^2}{24m\hbar}\sum_{j=1}^N\tau_j^3
=\frac{F}{2v\hbar}\sum_{j=1}^NS_j ,}\]

where \(S_j=vF\tau_j^3/(12m)\) is the parabolic segment cut off by the
\(j\)-th chord at horizontal speed \(v\), half of Lemma XI's tangent
area for that step. Three consequences follow.

\begin{itemize}
\item
  \textbf{Insertion.} Inserting a vertex at fraction \(\lambda\) of a
  step of duration \(\tau_j\) lowers the phase by
  \(F^2\tau_j^3\lambda(1-\lambda)/(8m\hbar)\), which is \(F/(2v\hbar)\)
  times the triangle the new vertex inscribes. At \(\lambda=\frac12\)
  this is Archimedes' triangle, three quarters of the segment, so
  repeated halving removes the phase in the proportions of his
  exhaustion, \(1:\frac14:\frac1{16}:\cdots\).
\item
  \textbf{Ordering.} The same Weyl increments composed in reverse order
  give the phase \(\Phi_N=F^2(\tau^3-\sum_j\tau_j^3)/(6m\hbar)\)
  relative to the chronological order, the largest over all orderings.
\item
  \textbf{Conservation.} For every partition,

  \[\Phi_N+4\theta_N=\frac{F^2\tau^3}{6m\hbar}=\frac{\tau\Delta E}{3\hbar}
  =\frac{F}{v\hbar}\,A_{\rm inertial,fall},\]

  so refinement moves phase from the polygon's segments into the
  ordering holonomy and keeps the total fixed at the Galileo area over
  \(\hbar\).
\end{itemize}

Both phases are c-numbers, so squeezing, mixing or any other choice of
body state leaves them unchanged. That is the property the two quantum
floors lacked, as corrected on 2026-09-22 in the
\href{mark-cost-and-statistical-floor.md}{mark-cost} and
\href{planck-gap-probabilistic.md}{probabilistic} notes. The single-step
case is the chord lens of the
\href{principia-constant-force-action.md}{constant-force note}, and the
closed loop of \href{newton-insertion-action.md}{N01} is another
instance of Lemma 2 below. New here are the identity for every
partition, the insertion law, the ordering holonomy and the conservation
law. Exploratory; no ledger promotion.

\hypertarget{setting-and-the-exact-lift}{%
\subsection{1. Setting and the exact
lift}\label{setting-and-the-exact-lift}}

Transverse coordinate \(\hat y\), momentum \(\hat p\),
\([\hat y,\hat p]=i\hbar\), mass \(m>0\); the horizontal motion at speed
\(v\) is common to every history considered and cancels. A force history
is an impulse density or measure \(f\) on \([0,\tau]\) with Hamiltonian
\(H_f(t)=\hat p^2/2m-f(t)\hat y\). In the interaction picture of the
free evolution \(U_0\), write \(W_f(t)=U_0(t)^\dagger U_f(t)\) and use
the Weyl operators of the probabilistic note,

\[D(a,b)=\exp\!\left[\frac{i}{\hbar}(b\hat y-a\hat p)\right],\qquad
D(z_1)D(z_2)=e^{-i\omega(z_1,z_2)/2\hbar}D(z_1+z_2),\]

with \(\omega(z_1,z_2)=a_1b_2-a_2b_1\). An impulse \(J\) at time \(t\)
is \(U_0(t)^\dagger e^{iJ\hat y/\hbar}U_0(t)=D(-Jt/m,J)\), since
\(U_0(t)^\dagger\hat yU_0(t)=\hat y+\hat pt/m\). The history's
\textbf{phase-plane path} is

\[z_f(t)=\bigl(\alpha_f(t),\beta_f(t)\bigr)
=\Bigl(-\frac1m\int_0^ts\,f(s)\,ds,\ \int_0^tf(s)\,ds\Bigr),\]

the accumulated transverse displacement relative to free motion and the
accumulated impulse. For the constant force, \(z_F(t)=(-Ft^2/2m,\ Ft)\),
a parabola in the phase plane.

\textbf{Lemma 1.}
\(W_f(\tau)=\exp\bigl[\frac{i}{2\hbar}\int_0^\tau\omega(z_f,dz_f)\bigr]\,D\bigl(z_f(\tau)\bigr)\).

\emph{Proof.} For finitely many impulses with increments
\(\delta z_1,\dots,\delta z_N\) in time order, induction on the Weyl
relation gives
\(D(\delta z_N)\cdots D(\delta z_1)=\exp[\frac{i}{2\hbar}\sum_{i<j}\omega(\delta z_i,\delta z_j)]D(\sum_j\delta z_j)\),
and
\(\sum_{i<j}\omega(\delta z_i,\delta z_j)=\sum_j\omega(z_{j-1},\delta z_j)\)
with \(z_{j-1}=\sum_{i<j}\delta z_i\). For a density \(f\) the generator
\(\frac{i}{\hbar}f(t)(\hat y+\hat pt/m)\) has c-number commutators at
different times, so the Magnus series stops at second order and the
Riemann sums converge to \(\int\omega(z_f,dz_f)\), as in N01 §4.
\(\square\)

\textbf{Lemma 2.} If \(z_f(\tau)=z_g(\tau)\), then
\(W_fW_g^\dagger=e^{i\mathcal A(f,g)/\hbar}\,\mathbf 1\), where
\(\mathcal A(f,g)=\frac12\oint\omega(z,dz)\) is the signed area enclosed
by \(z_f\) followed by \(z_g\) reversed.

\emph{Proof.} Lemma 1 for both, and Green's theorem for the closed loop.
\(\square\)

Two histories with equal endpoints leave the body in the same classical
state after \(\tau\), whatever happened in between, and Newton's
mechanics has no later observation that separates them. Lemma 2 says the
quantum evolutions still differ, by the enclosed phase-plane area over
\(\hbar\), and that the difference is a multiple of the identity. This
is the standard geometric phase of displaced oscillators, used in
trapped-ion phase gates
(\href{https://doi.org/10.1103/PhysRevA.62.022311}{Sørensen and Mølmer,
PRA \textbf{62}, 022311, 2000};
\href{https://doi.org/10.1038/nature01492}{Leibfried et al., Nature
\textbf{422}, 412, 2003}, both metadata). The \(t^3\) phase of a freely
falling wave packet is the single-history form of it
(\href{https://doi.org/10.1103/RevModPhys.51.43}{Greenberger and
Overhauser, RMP \textbf{51}, 43, 1979}, metadata).

\hypertarget{the-inscribed-polygon}{%
\subsection{2. The inscribed polygon}\label{the-inscribed-polygon}}

Partition \(0=t_0<\dots<t_N=\tau\) with \(\tau_j=t_j-t_{j-1}\) and
midpoints \(c_j=(t_{j-1}+t_j)/2\). \textbf{Newton's inscribed polygon}
gives the impulse \(F\tau_1/2\) at \(t_0\), \(F(\tau_j+\tau_{j+1})/2\)
at each interior \(t_j\), and \(F\tau_N/2\) at \(t_N\): each step's
impulse \(F\tau_j\) is split equally between its two ends.

\textbf{Theorem 1.} (i) On \((t_{j-1},t_j)\) the polygon moves uniformly
with momentum \(Fc_j\) along the chord of the parabola, and its
phase-plane path passes through \(z_F(t_j)\) at every vertex. (ii)
\(W_{\rm poly}(\tau)=e^{i\theta_N}W_F(\tau)\) with
\(\theta_N=F^2\sum_j\tau_j^3/(24m\hbar)\). (iii) With
\(S_j=vF\tau_j^3/(12m)\) the area between the \(j\)-th arc and its
chord, \(\theta_N=(F/2v\hbar)\sum_jS_j\).

\emph{Proof.} (i) After the second half of the impulse at \(t_{j-1}\)
the momentum is \(Ft_{j-1}+F\tau_j/2=Fc_j\), and in time \(\tau_j\) it
carries the body through \(Fc_j\tau_j/m=F(t_j^2-t_{j-1}^2)/2m\), the
parabola's transverse increment. In the phase plane the two
half-impulses of step \(j\) are \(k_1=(-F\tau_jt_{j-1}/2m,\ F\tau_j/2)\)
and \(k_2=(-F\tau_jt_j/2m,\ F\tau_j/2)\), whose sum is the parabola's
increment \(\delta z_j=(-F\tau_jc_j/m,\ F\tau_j)\).

\begin{enumerate}
\def\labelenumi{(\roman{enumi})}
\setcounter{enumi}{1}
\item
  Both evolutions are products over steps, in the same order, of
  operators with the same displacements \(\delta z_j\), so their ratio
  is the product of the per-step phase ratios. By Lemma 1 the polygon's
  step is \(D(k_2)D(k_1)=e^{i\omega(k_1,k_2)/2\hbar}D(\delta z_j)\) with
  \(\omega(k_1,k_2)=\frac{F\tau_j}{2}\cdot\frac{F\tau_j}{2m}(t_j-t_{j-1})=\frac{F^2\tau_j^3}{4m}\).
  The curve's step, with \(s=t-t_{j-1}\) and path
  \((-F(s^2+2st_{j-1})/2m,\ Fs)\) from \(z_F(t_{j-1})\), has
  \(\omega(z,dz)=\frac{F^2s^2}{2m}ds\), hence phase
  \(\frac1{2\hbar}\int_0^{\tau_j}\frac{F^2s^2}{2m}ds=\frac{F^2\tau_j^3}{12m\hbar}\).
  Neither depends on \(t_{j-1}\). The ratio is
  \(\exp[\frac{i}{\hbar}(\frac18-\frac1{12})\frac{F^2\tau_j^3}{m}]=\exp[\frac{iF^2\tau_j^3}{24m\hbar}]\).
\item
  Over the step, the chord exceeds the parabola by
  \(\frac{F}{2m}(t-t_{j-1})(t_j-t)\), and integrating against
  \(dx=v\,dt\) gives \(S_j=vF\tau_j^3/12m\). Then
  \(\frac{F}{2v}S_j=\frac{F^2\tau_j^3}{24m}\). \(\square\)
\end{enumerate}

Lemma XI's tangent area for the step is \(A_j=vF\tau_j^3/6m=2S_j\), so
\(\theta_N=(F/4v\hbar)\sum_jA_j\), and Proposition 1 of the
\href{planck-gap-paper.md}{paper} is the statement that this sum
vanishes with the mesh.

\textbf{Corollary 2 (insertion).} Replacing step \(j\) by steps
\(\lambda\tau_j\) and \((1-\lambda)\tau_j\) lowers \(\theta_N\) by

\[\frac{F^2\tau_j^3}{24m\hbar}\bigl[1-\lambda^3-(1-\lambda)^3\bigr]
=\frac{F^2\tau_j^3\,\lambda(1-\lambda)}{8m\hbar}
=\frac{F}{2v\hbar}\,T_j(\lambda),\]

with \(T_j(\lambda)=S_j-S_j'-S_j''\) the triangle inscribed between the
old chord and the new vertex. At \(\lambda=\frac12\),
\(T_j=\frac34S_j\), which is Archimedes' \emph{Quadrature of the
Parabola}, Propositions 17 and 24 (the segment is four thirds of its
inscribed triangle, by balancing and by exhaustion;
\href{../docs/classics/Archimedes_QuadratureParabola_Heath1897_OCR.md}{companion}).
After \(k\) rounds of halving every step, \(\theta\) is \(4^{-k}\) of
its value. \(\square\)

Each inserted point therefore changes the quantum evolution by a
computable phase, fixed by the area the insertion removes, and the
change is the same for every state of the body. A step's own phase
reaches one radian at

\[\tau_\hbar=\left(\frac{24m\hbar}{F^2}\right)^{1/3},\]

the scale that the constant-force note found for the single chord lens.
The insertion mesh of the mark theorem,
\(\tau_*=(48z_{1-\epsilon}^2m\hbar/F^2)^{1/3}\), has the same form,
obtained there from the cost of a mark and here from the phase of a
segment with no apparatus at all.

\hypertarget{the-ordering-holonomy-and-the-conservation-law}{%
\subsection{3. The ordering holonomy and the conservation
law}\label{the-ordering-holonomy-and-the-conservation-law}}

Keep the curve's step operators \(V_j=e^{i\varphi_j}D(\delta z_j)\) and
compose them in another order. For a permutation \(\sigma\), write
\(U_\sigma\) for the product with the steps applied in the order
\(\sigma\).

\textbf{Theorem 3.}
\(U_{\rm chron}U_\sigma^\dagger=\exp[\frac{i}{\hbar}\sum\omega(\delta z_i,\delta z_j)]\mathbf 1\),
the sum over the pairs \(i<j\) that \(\sigma\) inverts. Every term is
positive,

\[\omega(\delta z_i,\delta z_j)=\frac{F^2}{m}\tau_i\tau_j(c_j-c_i)>0\qquad(i<j),\]

so the reversed order is extremal, with

\[\Phi_N=\frac1\hbar\sum_{i<j}\omega(\delta z_i,\delta z_j)
=\frac{F^2}{6m\hbar}\Bigl(\tau^3-\sum_j\tau_j^3\Bigr).\]

\emph{Proof.} The per-step phases \(\varphi_j\) are common to every
ordering. By the Weyl relation the phase of a product is
\(\frac1{2\hbar}\) times the sum of \(\omega(\delta z_a,\delta z_b)\)
over pairs with \(a\) applied before \(b\); a pair kept in order cancels
between the two products, and an inverted pair contributes
\(\omega(\delta z_i,\delta z_j)-\omega(\delta z_j,\delta z_i)=2\omega(\delta z_i,\delta z_j)\).
With \(\delta z_i=(-F\tau_ic_i/m,\ F\tau_i)\),
\(\omega(\delta z_i,\delta z_j)=\frac{F^2}{m}\tau_i\tau_j(c_j-c_i)\).
For the sum,
\(c_j-c_i=\frac{\tau_i}2+\sum_{i<k<j}\tau_k+\frac{\tau_j}2\), so
\(\sum_{i<j}\tau_i\tau_j(c_j-c_i)=\frac12\sum_{i\ne j}\tau_i^2\tau_j+\sum_{i<k<j}\tau_i\tau_k\tau_j\),
and
\((\sum_j\tau_j)^3=\sum_j\tau_j^3+3\sum_{i\ne j}\tau_i^2\tau_j+6\sum_{i<k<j}\tau_i\tau_k\tau_j\)
gives \(\frac16(\tau^3-\sum_j\tau_j^3)\). \(\square\)

\textbf{Corollary 4 (conservation).}
\(\Phi_N+4\theta_N=F^2\tau^3/(6m\hbar)\) for every partition, and
\(F^2\tau^3/6m=\tau\Delta E/3=(F/v)A_{\rm inertial,fall}\) with
\(A_{\rm inertial,fall}=vF\tau^3/6m\) the paper's area (1). \(\square\)

In areas, \(\hbar\Phi_N=(F/v)(A-\sum_jA_j)\) and
\(4\hbar\theta_N=(F/v)\sum_jA_j\): the ordering holonomy carries the
part of the Galileo area that refinement keeps, the polygon defect
carries the part that Proposition 1 sends to zero, and the weighted
total is fixed. For the uniform partition,
\(\Phi_N=\frac{\tau\Delta E}{3\hbar}(1-N^{-2})\).

The chronological and reversed products realize the same displacement,
so the holonomy is a property of the composition law. Realizing a
reversed order in the laboratory needs displacements as well as
impulses, since an impulse applied at the wrong time must be translated
back; trapped-ion practice supplies such operations. Theorem 1 needs
only two force histories that act on a real body.

\hypertarget{what-the-phases-give-as-a-record}{%
\subsection{4. What the phases give as a
record}\label{what-the-phases-give-as-a-record}}

Put the two histories of Theorem 1 on the arms of a control qubit, as in
N01 §4. The body ends in the same state on both arms, so the qubit
acquires the relative phase \(\theta_N\) and nothing else, and the
optimal single-shot error in deciding polygon against curve is
\(\epsilon=\frac12(1-|\sin(\theta_N/2)|)\), with N01's formula for \(n\)
copies. A single step is resolved from its chord at error \(\epsilon\)
only if

\[\frac{F^2\tau_j^3}{24m\hbar}\ \ge\ 2\arcsin(1-2\epsilon),\qquad
\tau_j\ \ge\ \left(\frac{48\,m\hbar\,\arcsin(1-2\epsilon)}{F^2}\right)^{1/3},\]

for every state of the body, pure or mixed, squeezed or not. The bound
concerns this readout: the body's own displacement is a separate record,
priced by the mark and aperture theorems, and squeezing acts on that
record. What the phase supplies is the part of the signal that no
preparation can move.

\hypertarget{the-premise-restated}{%
\subsection{5. The premise, restated}\label{the-premise-restated}}

The identities above use the Weyl relations and nothing else: no state,
no variance, no apparatus. Of the two quantum premises used so far in
this programme, Robertson's inequality is the variance consequence of
the commutator, and the 2026-09-22 corrections show that squeezing moves
it. The commutator's phase content is what Theorems 1 and 3 use, and
squeezing leaves it fixed. So the premise that carries \(h>0\) into the
Galileo comparison without a resource bound is the \textbf{central
extension}: an impulse \(J\) and a displacement \(d\) compose, in
Newton's Corollary I, only up to the phase \(Jd/\hbar\) of their
parallelogram. Classically Corollary I is exact and commutative, the
polygon and the curve are the same state at every vertex, and every
\(\theta_N\) vanishes.

\textbf{Newton's fits supply a Newton-age form of that premise.} Two
passages of the \emph{Opticks}, now held verbatim in the
\href{../docs/classics/Newton_Opticks_1730_fits_and_queries.md}{companion},
fix how the interval of the fits varies with the corpuscle's speed. Book
II Part III Prop. XVII, ``manifest by the 10th Observation'', makes the
intervals in two mediums ``as the Sine of Incidence to the Sine of
Refraction'', so the interval is shorter in the denser medium by the
refractive ratio. Prop. X takes light to be ``swifter in Bodies than in
Vacuo, in the proportion of the Sines'', the hypothesis of Newton's
emission theory. Together they give

\[\Lambda\,v=\text{const across media},\qquad\text{hence}\qquad
\Lambda\,p=\text{const across media}\]

for each colour, the corpuscle's mass being unchanged by refraction.

\textbf{Proposition 5.} Suppose, extending Props. X and XVII from a step
in the medium to a continuously varying speed, that a corpuscle's fits
advance by one interval \(\Lambda=\Lambda_0p_0/|p|\) per length of path.
Then the number of fits along a path is its Maupertuis action
\(W=\int p\cdot dq\) in units of \(\Lambda_0p_0\). For the inscribed
polygon and the curve of Theorem 1 the counts differ by

\[\frac{\Delta W}{\Lambda_0p_0}=\frac{F^2}{12m\Lambda_0p_0}\sum_j\tau_j^3
=\frac{1}{2\Lambda_0p_0}\,\frac{F}{v}\sum_jA_j\]

fits, with \(A_j\) Lemma XI's tangent areas.

\emph{Proof.} Along the actual motion \(p\parallel dq\), so
\(|p|\,ds=p\cdot dq\) and the count is
\(\int p\cdot dq/(\Lambda_0p_0)\). The horizontal parts are common. On
step \(j\) the curve has
\(\int p_y\,dy=\frac{F^2}{m}\int_{t_{j-1}}^{t_j}t^2\,dt\) and the chord
has \(\frac{F^2}{m}c_j^2\tau_j\); their difference is
\(\frac{F^2}{m}\tau_j\bigl[\frac13(t_{j-1}^2+t_{j-1}t_j+t_j^2)-\frac14(t_{j-1}+t_j)^2\bigr] =\frac{F^2\tau_j^3}{12m}\),
and \(F^2\tau_j^3/6m=(F/v)A_j\). \(\square\)

In Newton's optics the fit at arrival decides whether the ray is
reflected or transmitted at the next surface (Prop. XII), so on his own
terms the polygon and the curve are optically different corpuscles once
their fits counts differ by half an interval, although they agree in
position and velocity at every vertex. Every ingredient is his: Lemma
XI's areas, the measured \(\Lambda=1/89000\) inch, the refraction
invariance of \(\Lambda p\) from Props. X and XVII, and the fits as a
determinate periodic disposition. The unknown is \(p\), which he could
not measure. Against the quantum statement, with \(\Lambda=\lambda/2\)
and \(\lambda p=2\pi\hbar\), the fits phase in radians is \(4\theta_N\):
one factor \(2\) because Newton's interval is half a wavelength, and one
because \(\Delta W=2\Delta S\) for this pair of histories.

The determinism of Prop. XII is untouched by this; what the fits supply
is a phase whose rate is proportional to momentum, and that is the
content of the central extension. So the premise that carries \(h>0\) in
the squeeze-immune form is present in Newton's optics, with
\(\Lambda p\) a positive refraction invariant. Two limits remain. Prop.
XV's rule for oblique emergence, a product of two secants, is more
complicated than a count of intervals along the path, so Proposition 5's
extension is ours. And \(\Lambda p\) is a per-colour constant in
Newton's system: universality across colours is what Planck's constant
adds.

\hypertarget{consequence-for-state}{%
\subsection{6. Consequence for STATE}\label{consequence-for-state}}

STATE's next item 1 is discharged in a stronger form than it was stated.
The state-independent quantity is identified for every refinement:
Newton's inscribed polygon differs from the parabola by
\(\theta_N=(F/2v\hbar)\sum_jS_j\), each inserted vertex changes it by
the inscribed triangle over \(\hbar\) in Archimedes' proportions, and
the ordering holonomy carries the complement, with
\(\Phi_N+4\theta_N=\tau\Delta E/3\hbar\). The Newton-age form of the
central extension is in §5: Props. X and XVII make \(\Lambda p\) a
refraction invariant, and the fits count separates the polygon from the
curve by \((F/v)\sum_jA_j/(2\Lambda p)\). Next: the general force law,
where Lemma 2 holds unchanged and Theorem 1's segment areas become the
areas between the phase-plane path and its chords; and a paper section
presenting Theorem 1 and Corollary 4 as the squeeze-immune core.
\end{document}
